\chapter{Introduction and Scenario}

\section{Scenario Choice}

This project follows \textbf{Scenario C - Omnichannel Commerce Core}. VerdeMart Retail
started by using nopCommerce as its online storefront, but the business now needs more than a
web shop. It wants web sales, warehouse execution, shipping, store operations, support, and
customer visibility to work together as part of one commerce experience.

In this scenario, nopCommerce becomes the commerce core of that wider ecosystem. A web
order should be reflected in operational systems such as warehouse and shipping, and changes
that happen outside nopCommerce, such as stock or fulfillment updates, should become visible
again in the customer and operator experience. This is important because those surrounding
systems may be slow, stale, unavailable, or inconsistent, but the commerce experience still
needs to remain useful and understandable.

\section{Core Architectural Problem}

Currently, nopCommerce mainly acts as the online shop: it owns the storefront, catalog
browsing, cart, checkout, customer accounts, and order creation inside one web-centered
platform. This works while the business is mostly operating through the online channel, but it
does not fully address an omnichannel environment where warehouse execution, shipping,
store operations, and support systems also influence the customer experience.

The core architectural problem is deciding how this online shop can evolve into a commerce
core without assuming that every surrounding operational system is always fast, available, and
consistent. The architecture must define what remains inside nopCommerce, what moves around
it, and how the system behaves when warehouse or shipping state is delayed, stale, unavailable,
or later recovered.
